\documentclass[conference]{IEEEtran}
\IEEEoverridecommandlockouts
% ────────────────────────────────────────────────────────────────────────
% DRAFT generated from project README + session logs. Every \cite{} key is
% a PLACEHOLDER -- see the "References (to fill in)" comment block at the
% end of this file for what each key needs and what is already known about
% it. Add the real BibTeX entries to refs.bib yourself before compiling.
% All tables/numbers are taken directly from this project's own recorded
% experiment results (README.md and this session's eval batches), not
% fabricated.
% ────────────────────────────────────────────────────────────────────────
\usepackage{cite}
\usepackage{amsmath,amssymb,amsfonts}
\usepackage{algorithmic}
\usepackage{graphicx}
\usepackage{textcomp}
\usepackage{xcolor}
\usepackage{booktabs}
\usepackage{multirow}
\usepackage{url}

\def\BibTeX{{\rm B\kern-.05em{\sc i\kern-.025em b}\kern-.08em
    T\kern-.1667em\lower.7ex\hbox{E}\kern-.125emX}}

\begin{document}

\title{Soft Body Identification from Optical Flow: \\
Unsupervised Instrument and Soft-Tissue Segmentation \\
in Surgical Video via Motion Reconstruction}

\author{
\IEEEauthorblockN{1\textsuperscript{st} Tianming Zhang}
\IEEEauthorblockA{}
}

\maketitle

\begin{abstract}
Pixel-level annotation of surgical video is expensive and slow to
collect, which limits fully-supervised instrument and tissue
segmentation to whichever datasets happen to be annotated and leaves
such models generalizing poorly across hospitals, scopes, and
procedures. We address this with an unsupervised approach that uses
motion, rather than manual labels, as the supervisory signal for both
instrument and soft-tissue segmentation. Building on RCF
(Rigid-Cluster-Flow)~\cite{rcf_cvpr2023}, which reframes segmentation as
motion clustering -- a soft per-pixel cluster assignment is learned by
requiring each cluster's own rigid motion model to reconstruct optical
flow from a frozen RAFT~\cite{raft_eccv2020} network -- we extend the
method to a new, heterogeneous surgical dataset (``CMC''), for
instrument segmentation and, novel to this work, for soft-tissue
segmentation, a target whose motion signal is an order of magnitude
weaker and noisier than the instrument's. We add a DINO-feature
consistency loss~\cite{dino_iccv2021} to stabilize which output channel
represents which object across training, and report roughly 150
controlled single-variable experiments spanning loss design,
segmentation-head architecture, training/data-quality interventions,
multi-frame extensions, and a DINO graph-partitioning
prior~\cite{graphpartition_mchm2024}. Across this arc we distill one
recurring, reproducible finding: gains come almost exclusively from
protecting and denoising an already-scarce motion signal, not from added
model capacity. We close with a diagnosis of the project's current
performance ceiling and a concrete proposal -- explicit, filtering-style
motion decoupling with a dedicated sub-network, in the spirit of recent
dynamic scene reconstruction work~\cite{batrack2025} -- for what would
need to change to move past it.
\end{abstract}

\begin{IEEEkeywords}
unsupervised segmentation, optical flow, surgical video, motion clustering, self-supervised learning
\end{IEEEkeywords}

\section{Introduction}

Minimally invasive surgery relies heavily on video for both intra-operative
guidance and post-operative analysis, and automatic segmentation of
surgical instruments and soft tissue is a prerequisite for most
downstream applications (workflow analysis, skill assessment, augmented
reality overlays). Fully-supervised segmentation models are the dominant
approach, but they require per-frame pixel masks that are costly to
collect and that do not transfer well across hospitals, scopes, and
procedures with different visual statistics. This motivates
\emph{unsupervised} approaches that use a signal already present in the
raw video -- motion -- instead of manual labels.

RCF~\cite{rcf_cvpr2023} formulates unsupervised video object segmentation
as an Expectation-Maximization-like motion clustering problem implemented
through regression: a convolutional network predicts a soft, per-pixel
cluster assignment (the segmentation mask); each cluster's own rigid (or
near-rigid) motion model is then fit, in closed form, to the optical flow
under that cluster's current soft assignment (the ``M-step''); and the
whole pipeline is trained end-to-end by requiring the mask-weighted
mixture of per-cluster motion predictions to reconstruct the externally
computed, frozen optical flow (RAFT~\cite{raft_eccv2020}) between adjacent
frames. No manual segmentation label ever enters training -- only the
optical flow, an external, non-trainable signal, does.

This project starts from that idea and pursues it in two stages. First,
because RCF's own paper does not release a pretrained checkpoint or
training configuration for the instrument-segmentation setting, we
reproduce it from scratch on the public EndoVis~2017 instrument
segmentation benchmark~\cite{endovis2017}, and find that a
straightforward reproduction is unstable: the instrument does not
reliably stay bound to one output channel across training. We add an
auxiliary DINO-feature consistency loss~\cite{dino_iccv2021} to anchor the
instrument channel by appearance, which resolves the instability and
becomes the standard configuration for everything that follows. Second,
we transfer this recipe to a private, considerably more heterogeneous
laparoscopic dataset (``CMC'', collected across multiple grasp-event
offsets and, later, multiple hospitals/resolutions), first for instrument
segmentation and then -- on top of an instrument-competent model -- for
soft-tissue segmentation, where the motion signal is an order of
magnitude weaker (soft tissue moves less, and less rigidly, than a
held instrument) and correspondingly harder to learn from without
collapsing.

Across roughly 150 numbered, single-variable experiments recorded in this
project's history, we find a consistent pattern: architectural capacity
increases (deeper segmentation heads, ASPP, U-Net-style decoders, extra
injected features, multi-frame fusion) plateau quickly, while the only
mechanism-level changes that produced clean, reproducible improvements
were ones that protected or denoised the existing motion-reconstruction
signal (supervision density around motion boundaries, removing corrupted
flow pairs, correcting a downsampling operation that blended
object-boundary motion into the background before the loss ever saw it).
This paper documents that arc, the specific architecture and training
protocol at each stage, the ablations that did and did not work, and a
concrete diagnosis -- backed by a within-project literature detour into
recent bundle-adjustment-based dynamic scene reconstruction
work~\cite{batrack2025} and a re-derivation of a DINO
graph-partitioning eigenvector prior~\cite{graphpartition_mchm2024} in
explicit Bayesian terms -- of what is currently limiting further
progress, and what a principled next step looks like.

\textbf{Contributions.}
\begin{itemize}
\item A from-scratch reproduction of RCF's instrument-segmentation
setting on EndoVis~2017, with an added DINO consistency loss that fixes a
channel-instability failure mode absent from the original paper's
description, reaching approximately 70\% mIoU.
\item A transfer of this recipe to a new, heterogeneous surgical dataset
(CMC) for both instrument and (novel to this project) soft-tissue
segmentation, including a diagnosis and fix for a resolution/field-of-view
domain gap between the source and target datasets.
\item A systematic experimental study (data-quality/signal-density
interventions, segmentation-head architecture search, multi-frame
extensions, DINO graph-partitioning fusion) isolating which classes of
change reliably help this specific unsupervised motion-reconstruction
architecture and which do not.
\item A root-cause diagnosis of a previously unexplained instrument-mIoU
decay observed in the strongest recent configuration, traced to a
specific background-motion-removal mechanism's lack of robustness to
foreground (instrument) contamination.
\item A reformulation of DINO graph-partitioning eigenvector fusion as an
explicit Bayesian prior (rather than an opaque mid-network feature
concatenation), combined with an existing-but-unused motion-derived
likelihood term already implicit in the architecture's own EM structure.
\end{itemize}

\section{Related Work}

\textbf{Unsupervised motion-based video object segmentation.} This
project builds directly on RCF~\cite{rcf_cvpr2023}, which frames
segmentation as regression-based EM motion clustering driven by optical
flow reconstruction. The core signal is a frozen, pre-trained optical
flow estimator, RAFT~\cite{raft_eccv2020}, applied off-the-shelf as an
external target -- the network never receives motion supervision beyond
what RAFT itself outputs, so the ceiling on segmentation quality is
coupled to RAFT's own accuracy on the target domain.

\textbf{Self-supervised appearance features.} DINO~\cite{dino_iccv2021}
vision transformers, trained with no labels via self-distillation, are
used twice in this project: (i) as an auxiliary per-pixel consistency
loss that anchors channel identity during training (Section
\ref{sec:endovis}), and (ii) as a frozen feature extractor for graph-based
semantic partitioning (Section \ref{sec:dino_graph}).

\textbf{Graph-partitioning segmentation from self-supervised features.}
\cite{graphpartition_mchm2024} shows that a cosine-similarity affinity
graph built from frozen DINO patch features, spectrally decomposed via
the normalized graph Laplacian~\cite{normalized_cuts}, yields eigenvectors
that meaningfully partition a surgical scene into semantic modules
(instrument, tissue, background) with no supervision at all -- entirely
independent of motion. We revisit this mechanism in Section
\ref{sec:dino_graph} and re-derive it as a genuine Bayesian prior over the
same cluster assignment RCF's motion likelihood already estimates,
rather than treating it as an opaque feature to concatenate.

\textbf{Dynamic scene reconstruction via explicit motion decoupling.}
\cite{batrack2025} decouples camera-induced (static) motion from
object-induced (dynamic) motion using a dedicated, separately-trained
network, then hands the isolated static component to classical bundle
adjustment. Their ablation shows that (a) explicit decoupling and (b) a
learned camera/dynamic-region mask compose multiplicatively (roughly a
$4\times$ error reduction together, each individually contributing a
factor of $\sim2\times$), and, notably, that training one network to
predict \emph{both} the total and the dynamic motion component is
worse than two separate networks. We return to this finding in Section
\ref{sec:future} as the basis for a concrete architectural proposal.

\section{Method}

\subsection{Base architecture: motion reconstruction as regression-EM}
\label{sec:base_arch}

The network has a shared ResNet-50 backbone~\cite{resnet_cvpr2016}
(dilated in the last two stages to preserve spatial resolution),
initialized from DenseCL self-supervised pretraining, feeding three
heads:

\begin{itemize}
\item \textbf{Mask head} (multi-scale segmentation head with an ASPP
module~\cite{aspp_deeplab}, Section~\ref{sec:archsearch}): produces a
$C$-channel softmax over the backbone's fused multi-scale features -- the
soft cluster assignment (E-step).
\item \textbf{Motion head}: for each of the $C$ channels, fits a
per-channel rigid motion model (constant translation, or translation with
an affine correction) to the optical flow using the current mask as a
weighted-least-squares weight (M-step), then adds a learned, bounded
per-pixel residual to absorb non-rigid deformation the rigid model cannot
explain.
\item \textbf{Residual head}: an independent branch predicting a
secondary per-channel residual correction directly from the deepest
backbone feature, decoupled from the mask.
\end{itemize}

The mask-weighted mixture of the $C$ channels' motion predictions is
compared against the RAFT flow with an $L_2$ reconstruction loss, computed
\emph{only} near detected motion boundaries (an angle-change detector on
the flow field, dilated and floored so non-boundary pixels retain a small
residual weight) -- this supervision-density mechanism, introduced
partway through the project (Section~\ref{sec:signal_quality}), turned
out to be the single most consequential change in the whole history of
the project. Per training step, an easy-example-mining rule keeps only
the $k$ lowest-loss samples in the batch for backpropagation.

No manual segmentation label participates in this loss. Channel identity
(which of the $C$ channels ends up representing the instrument, which
represents tissue, etc.) is not fixed by architecture; it emerges from
training and is resolved at evaluation time by an oracle matching
procedure (Section~\ref{sec:eval_protocol}).

\subsection{EndoVis~2017 reproduction and the DINO consistency loss}
\label{sec:endovis}

RCF's own paper does not release pretrained weights or a training
configuration for surgical instrument segmentation, so this project first
reproduces the method from scratch on EndoVis~2017~\cite{endovis2017}.
A direct reproduction of the base architecture (Section~\ref{sec:base_arch})
is unstable at this stage: because the loss only ever supervises the
\emph{mixed} reconstruction, not any individual channel's identity, the
instrument's soft assignment can drift across channels over training,
and different channels can partially split or merge the instrument
region with no signal telling the network this is wrong.

We add an auxiliary DINO-feature consistency loss: for each channel $c$,
minimize the mask-weighted mean cosine distance between every assigned
pixel's frozen DINO feature and that channel's own mask-weighted mean DINO
feature,
\[
\mathcal{L}_{\text{dino}} = \frac{1}{C}\sum_{c}
\frac{\sum_p \hat M_c(p)\,\big(1-\cos(f_p,\mu_c)\big)}{\sum_p \hat M_c(p)},
\]
where $\mu_c$ is the mask-weighted mean of $\ell_2$-normalized DINO
features under channel $c$. This does not require or assume any specific
channel index for the instrument -- it only pushes each channel toward
visually coherent membership, which in practice is enough to keep the
instrument bound to a stable identity across training. This loss (small
weight, $w_{\text{dino}} \approx 0.05$--$0.1$) becomes the standard
configuration for every later experiment in this project. With it, the
EndoVis reproduction reaches approximately 70\% instrument mIoU,
establishing the base recipe subsequent sections build on.

\subsection{Transfer to CMC: instrument segmentation}
\label{sec:cmc_instrument}

Directly fine-tuning the EndoVis-trained model on the CMC dataset
under-performed. CMC is collected across multiple source hospitals with
different camera resolutions and framing conventions, and the
EndoVis-trained backbone carries an inductive bias (in particular, a
learned notion of object scale/field-of-view) that does not transfer
cleanly. Rather than discarding the EndoVis prior outright, we diagnosed
the dominant failure mode as a scale/resolution mismatch (the network's
effective receptive field vs. object scale, inherited from EndoVis's own
image statistics) and adjusted the resize/crop configuration used during
fine-tuning so that the CMC crop presented to the network matches the
scale distribution the backbone was pretrained under, closing most of the
gap. With this fix, fine-tuning on CMC reaches roughly 67\% instrument
mIoU -- below the EndoVis number, attributable to (a) the CMC dataset
being considerably smaller, and (b) the well-known EM/regression tension
between a strong pretrained prior and the (weaker, noisier) new evidence
the posterior update has to reconcile it with. Sliding-window inference
(Section~\ref{sec:sliding_window}) is introduced at this stage to handle
CMC's variable, generally larger image resolutions without retraining at
native resolution.

\subsection{Soft-tissue segmentation: loss-engineering attempts (largely unsuccessful)}
\label{sec:tissue_loss}

With an instrument-competent CMC model as a starting point, we next
target soft-tissue segmentation by continuing training on grasp-event
data with additional supervision meant to shape a dedicated tissue
channel while protecting the already-learned instrument channel. Roughly
20 auxiliary losses were tried across this phase, targeting: rigidity
(reward high intra-channel flow variance, i.e. non-rigid deformation, for
candidate tissue channels), motion divergence (grasped tissue compresses
under contact, producing non-zero flow divergence RAFT can detect),
grasp-point convergence, cross-channel flow-direction clustering (both
self-referential, using the model's own current mask as the pseudo-label,
and externally referenced, using an independent $k$-means clustering of
RAFT flow direction), and several distillation variants (one-sided vs.\
symmetric BCE) meant to protect the instrument channel from being
absorbed by the new tissue-oriented pressure.

None of these auxiliary losses produced a stable, above-baseline result:
runs oscillate close to the instrument-only baseline, and several
specific, diagnosable failure modes recur -- most commonly the instrument
channel degrading over training as increasingly strong tissue-oriented
pressure (deformation reward, entropy regularization) wins the underlying
softmax competition for pixels near the instrument/tissue boundary. By
the end of this phase every one of these losses is disabled (weight 0)
in the surviving configuration; the two losses that remain active
throughout the rest of the project are the core motion-reconstruction
loss and the small DINO consistency loss from Section~\ref{sec:endovis}.
We conclude that the inductive bias carried over from the instrument-only
pretraining was not something an additional loss term could cheaply
redirect, and that architecture, not loss design, was the more promising
lever -- motivating Section~\ref{sec:archsearch}.

\subsection{Segmentation-head architecture search}
\label{sec:archsearch}

Observing that the segmentation head only consumed two of the backbone's
four available multi-scale feature maps, we trained the tissue model from
scratch on the merged grasp0+grasp10 CMC data with a series of
segmentation-head architectures: a U-Net-style top-down decoder with skip
connections, several ASPP/multi-scale-fusion variants, DenseASPP
(cascaded dilated convolutions), and strip pooling (targeting the
instrument's characteristically elongated shape). Table~\ref{tab:arch}
summarizes the two configurations that mattered most.

\begin{table}[htbp]
\caption{Segmentation-head architecture search, selected results (mIoU \%, oracle-matched)}
\label{tab:arch}
\centering
\begin{tabular}{lccc}
\toprule
Version & Change & Instrument & Tissue \\
\midrule
v64 & MultiScaleSegHead + ASPP (dil.\ 6/12/18 + GAP) & 65.57 & 73.32 \\
\bottomrule
\end{tabular}
\end{table}

MultiScaleSegHead with a parallel-dilation ASPP module, replacing a
single fusion convolution, was the only architecture change in this
search that produced a clean, reproducible improvement (v64, sum
mIoU 138.89, the project's best result up to that point). Subsequent
refinement attempts on top of this head (finer dilation schedules,
additional pooling branches, alternative skip-connection patterns)
consistently failed to exceed v64, indicating the segmentation head's own
capacity had saturated for this data scale -- further gains, documented
next, came from the training signal itself rather than the network.

\subsection{Signal density and data quality}
\label{sec:signal_quality}

Two changes, made independently of architecture, produced this project's
largest and most reproducible gains, and are the basis for this paper's
central empirical claim.

\textbf{Boundary supervision density (v83).} The reconstruction loss's
boundary mask (Section~\ref{sec:base_arch}) was, at this point, a thin,
un-dilated band around detected motion-angle discontinuities, covering
only $\sim$20\% of pixels; every other pixel received zero loss weight.
Widening this band (dilation $7\!\to\!15$~px) and giving the remaining,
non-boundary pixels a small non-zero floor weight (rather than exactly
zero) raised effective supervision coverage to $\sim$40\% of pixels.
Combined with removing 24 training pairs whose RAFT flow had failed
outright (identified by extreme flow magnitude, plausible causes:
motion blur or textureless/specular frames RAFT cannot track), this
produced sum mIoU 139.36 (instrument 66.76, tissue 72.60) -- a new record
and, per this project's own repeated later re-verification, the only
single mechanism-level change in the project's entire history that
reliably improved on its predecessor whenever it was combined with other
new mechanisms.

\textbf{Downsampling mode for the ground-truth flow (v102).} The
ground-truth flow is resized from crop resolution ($384\!\times\!384$) to
the mask resolution ($96\!\times\!96$) before the reconstruction loss is
computed. The default (bilinear) interpolation blends instrument-boundary
motion vectors with surrounding background motion \emph{at the
downsampling step itself} -- for a target occupying roughly 1\% of frame
area, this directly works against the density mechanism just described,
diluting exactly the pixels it was built to protect. Switching to area
(box) downsampling, the more standard choice for \emph{downsampling}
specifically, gave 139.77 (0.30 threshold) / 140.12 (0.35 threshold) --
this project's champion result prior to this session's work, and the
baseline every subsequent experiment (Sections
\ref{sec:multiframe}--\ref{sec:this_session}) is compared against.

We highlight this pattern because of how it recurs: two of the three
best-performing changes across $\sim$150 experiments are corrections to
how an already-present signal was being computed or supervised, not new
information sources.

\subsection{Multi-frame extensions}
\label{sec:multiframe}

All experiments up to this point predict the mask independently per
frame, with no mechanism for one frame's prediction to be informed by
neighboring frames. We tested whether giving the mask branch access to
multiple frames -- either via an approximate ``bridge'' construction
(reusing existing single-gap frame pairs across grasp-event offsets to
simulate a larger temporal gap) or via genuine multi-frame continuous
video later obtained for a subset of cases -- would help. Concretely:
joint cross-frame mask architectures (shared features across a frame
pair, feature concatenation, a 4-scale extension, and a deformable
cross-frame attention variant) were each tried on top of the v102
baseline. None produced a clean improvement; several of the largest,
most structurally novel variants (deformable attention in particular)
were the weakest. We separately confirmed, and fixed, a real
long-standing bug discovered during this arc: the project's own
random-gap sampling mechanism, present in the dataset code since early in
the project, had been silently dead for every CMC configuration because
CMC's directory structure gave every non-unit gap draw the same source
folder as the gap-1 draw, causing it to fall back to gap-1 with no
warning. We conclude, tentatively, that the current mask-prediction
architecture -- and, per Section~\ref{sec:signal_quality}'s pattern, the
current bottleneck more broadly -- is not well suited to naive multi-frame
feature fusion; modeling genuine cross-frame relationships appears to
need a different mechanism than concatenation or attention over
independently-extracted per-frame features (see Section~\ref{sec:future}).

\subsection{DINO graph-partitioning eigenvector fusion}
\label{sec:dino_graph}

\cite{graphpartition_mchm2024} shows that eigenvectors of the normalized
graph Laplacian of a cosine-similarity affinity graph over frozen DINO
patch features meaningfully partition a surgical frame into semantic
modules with no supervision, no motion, and no training at all -- we
verified this reproduces cleanly on our own frames using a frozen DINO
ViT-Small/8 at $384\!\times\!384$ input, downsampled to a $32\!\times\!32$
graph, and that the identical affinity/Laplacian/eigendecomposition
pipeline run on our own trained ResNet-50 features instead produces pure
noise -- DINO, not the project's own backbone, is required as the feature
source for this to work.

We tested seven independent fusion mechanisms/injection points on top of
the v102 baseline, spanning: concatenating eigenvectors into different
backbone feature stages, a non-parametric self-referential centroid
matching mechanism, injecting eigenvector-derived confidence directly
into the motion-fitting step's weights, a multiplicative attention gate,
simultaneous multi-scale injection, and decoder-stage injection after the
segmentation head's own multi-scale fusion. Table~\ref{tab:dino_batch}
summarizes.

\begin{table}[htbp]
\caption{DINO graph-partitioning fusion, seven mechanisms vs.\ v102 baseline}
\label{tab:dino_batch}
\centering
\begin{tabular}{lccc}
\toprule
Version & Mechanism & Val.\ sum peak & Test mIoU \\
\midrule
v102 (baseline) & -- & 139.77 & -- \\
v146 & concat $\to$ feat0 (shallow) & 137.0 & 73.55 \\
v149 & concat $\to$ feat3 (deep) & \textbf{139.0} & \textbf{74.57} \\
v147 & self-ref.\ centroid matching & 137.0 & 69.71 \\
v148 & inject into motion-fit weights & 133.0 & 67.52 \\
v153 & multiplicative attention gate & 126.0 & 73.37 \\
v154 & feat2 $+$ feat3 (double inject.) & 131.0 & 68.28 \\
v155 & decoder-stage injection & 113.0 & 67.11 \\
\bottomrule
\end{tabular}
\end{table}

The simplest mechanism (concatenation into the deepest, most semantic
backbone feature stage, v149) was the best of the seven -- every
mechanism explicitly designed to be ``smarter'' than plain concatenation
(gating, non-parametric matching, direct injection into the motion fit)
underperformed it, several substantially. We interpret this as the
opposite of a coincidence: a fixed-form combination rule imposes a prior
on \emph{how} two independent signal sources should be combined, and
under this project's small data regime (on the order of 1500 training
pairs), that prior is more often wrong than the flexibility a plain
learned convolution retains.

Independently, we tested a more robust background-motion removal step
(subtracting an estimated camera/background flow, fit via
Iteratively-Reweighted Least Squares, before the per-channel motion fit)
on top of v149. The plain (non-robustified) variant, v152, reached the
highest single validation peak recorded in this batch (140.0), matching
the v102 champion range, but a full-training-curve analysis revealed a
systematic, previously unnoticed problem, detailed next.

\subsection{Root-cause diagnosis: background-motion removal and instrument decay}
\label{sec:bg_removal}

Pulling the full, per-epoch validation curve (not just the sparse
checkpoints normally used for comparison) for four closely related
configurations reveals a clear pattern invisible in checkpoint-level
comparisons alone: configurations without background-motion removal
(v102, v149) show noisy but flat instrument mIoU across all 80 training
epochs, while both variants \emph{with} background-motion removal (v151,
IRLS-robustified; v152, plain least-squares) show a distinct
\emph{early-peak-then-decline} shape -- a strong result in the first
$\sim$20 epochs, followed by a systematic decline to a $\sim$15-point
lower plateau for the remainder of training. This isolates the
background-removal mechanism itself, independent of the DINO fusion it
was layered on top of, as the cause.

The root cause, stated in that mechanism's own pre-existing
implementation notes but never previously connected to an observed
failure mode, is that the background-motion fit is not robust to
instrument pixels biasing it. Critically, the residual-magnitude-based
robustification (IRLS) that was tried does not fix this: a rigidly
translating instrument has smooth \emph{internal} flow, so it produces no
local residual spike relative to its own motion anywhere except at its
own silhouette edge -- IRLS can only ever detect and downweight the
edge, never the interior, and so under-corrects exactly the region that
matters. An attempted fix using an EMA teacher's own instrument-channel
prediction as a semantic (not residual-magnitude-based) exclusion weight
for this fit was tested and found, empirically, to underperform both the
uncorrected baseline and the original mechanism, plausibly because the
EMA teacher's segmentation head is itself too unreliable early in
training to serve as a dependable exclusion signal, even with a
warm-up delay.

\subsection{Train/eval scale mismatch, multi-scale training, and this session's remaining changes}
\label{sec:this_session}

A separate, previously undiagnosed issue was found in the evaluation
pipeline: training's data augmentation resizes source images differently
depending on their native resolution (two discrete target scales), while
evaluation resizes every image to one universal target scale regardless
of source resolution. For the higher-resolution source images specifically,
this amounts to a real $\sim\!1.47\times$ scale mismatch between what the
network was trained to expect and what it is shown at evaluation time
(objects appear noticeably smaller at eval than during training).
Patching evaluation alone to match training's resize logic, without
retraining, made results \emph{worse}, not better (140.19 $\to$ 133.16),
because the champion checkpoint's own selection (which checkpoint counted
as ``best'' during training) had itself always been computed under the
old, mismatched evaluation scale -- correcting evaluation alone breaks
that calibration rather than fixing anything. This motivated fixing the
training side instead: replacing the two discrete resize targets with one
continuous, source-resolution-independent random-scale range (standard
multi-scale training), training for more epochs to compensate for the
wider distribution needing more iterations to converge, and (separately)
enabling this project's existing-but-previously-unused exponential moving
average (EMA) weight-smoothing mechanism to reduce the training curve's
substantial epoch-to-epoch variance. We also reformulated the DINO
graph-partitioning fusion of Section~\ref{sec:dino_graph} as an explicit
Bayesian prior (Section~\ref{sec:bayesian_prior}) rather than a
feature-level concatenation, combined with the project's own,
previously-unused motion-derived likelihood term. Results for this batch
of changes are reported alongside the rest in
Section~\ref{sec:results_summary}.

\subsection{DINO eigenvectors as a Bayesian prior}
\label{sec:bayesian_prior}

Re-reading~\cite{graphpartition_mchm2024} in full (rather than only its
affinity/Laplacian recipe) surfaces an important correction to how this
project had informally described the eigenvector signal: an individual
eigenvector is not a per-pixel class probability. Each is a single
real-valued partition ``mode'' at a specific normalized-cut granularity
-- the lowest-eigenvalue eigenvectors give the coarsest, most salient
split (typically instrument vs.\ background), higher-index eigenvectors
give progressively finer sub-structure, and eigenvectors beyond roughly
the 10th--15th become noise-like. The paper's own recipe for turning
this into an actual segmentation is to stack the leading $g$ eigenvectors
into a $g$-dimensional per-pixel embedding and run $k$-means on that
embedding -- not to read class membership off the raw values.

This project's own architecture (Section~\ref{sec:base_arch}) is already,
explicitly, a regression form of Expectation-Maximization: the mask is a
soft E-step responsibility, and the per-channel motion fit is an M-step.
Under that framing, the DINO eigenvector embedding is a genuine,
motion-independent \emph{prior} $P(\text{channel}\mid\text{appearance})$,
and the architecture already contains, unused, a genuine motion-derived
\emph{likelihood}: an auxiliary consistency term whose (previously never
enabled) target is
\[
P(\text{channel}=c \mid \text{flow}) \;\propto\;
\exp\!\Big(-\tfrac{1}{T}\,\big\lVert
\hat F_c(\text{rigid model}) - F_{\text{RAFT}} \big\rVert^2\Big),
\]
i.e.\ how well channel $c$'s \emph{own} fitted rigid motion model
explains the observed flow at each pixel. Bayes' rule combines a prior
and a likelihood by adding their logs. We implement this directly:
$C$ learnable centroids in the $g$-dimensional eigenvector-embedding
space give a temperature-scaled soft-distance prior logit, of the same
mathematical form as the motion likelihood above, added straight onto the
segmentation head's raw mask logits before the final softmax -- a
decision-layer fusion, in contrast to Section~\ref{sec:dino_graph}'s
feature-layer concatenation. The added module is initialized as an exact
no-op (a zero-initialized output scale), so training begins at the
baseline's own optimum. We additionally corrected a design inconsistency
found while implementing this: the motion-likelihood auxiliary loss, as
originally written, applied to every pixel uniformly, while the
project's main reconstruction loss is deliberately concentrated near
motion boundaries (Section~\ref{sec:signal_quality}) -- we made the two
losses agree on where the signal should matter by having the auxiliary
loss reuse the same boundary weighting.

\section{Evaluation Protocol}
\label{sec:eval_protocol}

\subsection{Oracle matching and channel identity}
Because channel identity is not fixed by the architecture, evaluation
must resolve, per checkpoint, which output channel(s) correspond to the
instrument and which to tissue. Two protocols are used across this
project, and are not equivalent: (i) an \emph{oracle} protocol, used for
per-epoch training-time monitoring, that exhaustively searches every
non-empty subset of channels per frame independently and reports the
best-matching union's IoU -- an optimistic, best-case-given-this-prediction
number used only for checkpoint selection; and (ii) a \emph{locked}
protocol, used for final reported instrument numbers, that fixes a
channel (or channel pair) once, from a held-out detection pass, and
applies it uniformly. Tissue numbers reported throughout this paper
remain oracle (channel(s) already known to be instrument are excluded
from the candidate set, but the remaining channels are still searched
exhaustively per frame) -- a caveat we flag explicitly, since it means
reported tissue numbers are somewhat more optimistic than reported
instrument numbers, and this asymmetry should be kept in mind when
comparing across versions.

\subsection{Sliding-window inference}
\label{sec:sliding_window}
Because the network is trained on fixed $384\!\times\!384$ crops but CMC
images are larger and of variable resolution, test-time inference uses a
sliding window ($384$px, $192$px stride) with probability-space averaging
in overlapping regions. We tested two alternatives: (i) hard
nearest-window-center assignment (no blending across windows, in the
spirit of the original U-Net's overlap-tile strategy), and (ii) a
distance-weighted (Hann-window) taper. Both alternatives underperformed
the existing probability-averaging default (by $-1.58$ and $-0.33$ to
$-1.09$ points respectively, depending on threshold), most of the loss
concentrated on the tissue channel, whose boundaries are more diffuse
than the instrument's -- we conclude the current default should be kept.

\section{Results Summary}
\label{sec:results_summary}

\begin{table}[htbp]
\caption{Top results across the full v1--v134 project history (sum mIoU, oracle-matched)}
\label{tab:top10}
\centering
\begin{tabular}{clc}
\toprule
Rank & Version & Sum mIoU \\
\midrule
1 & v102 & 139.77 (140.12 at threshold 0.30) \\
2 & v83  & 139.36 \\
3 & v64  & 138.89 \\
4 & v75  & 138.50 \\
5 & v47  & 138.28 \\
6 & v95  & 137.89 \\
7 & v43  & 137.72 \\
8 & v85  & 137.59 \\
9 & v91  & 137.59 \\
10 & v63 & 137.57 \\
\bottomrule
\end{tabular}
\end{table}

Table~\ref{tab:top10} lists the ten best results across the first 134
versions of the project (the multi-frame and DINO-graph lines, discussed
above, are evaluated separately since they were not directly comparable
under the same protocol at the time). The top three results each
represent a distinct class of intervention -- a data-preprocessing
correction (v102), a supervision-density and data-cleaning change (v83),
and an architecture change (v64) -- and, consistent with the pattern
argued throughout this paper, the two highest-ranked results are both
signal-quality corrections rather than new architecture or new
information sources.

\begin{table}[htbp]
\caption{This session's changes vs.\ the v102/v152 baselines, best checkpoint per configuration (grasp0 + grasp10 combined, sum mIoU)}
\label{tab:this_session}
\centering
\begin{tabular}{llc}
\toprule
Version & Change vs.\ base & Combined sum \\
\midrule
v102 (base) & -- & 282.85 \\
v152 (base) & DINO fusion + plain bg-removal & \textbf{282.89} \\
v151 & \quad + IRLS-robustified bg-removal & 281.82 \\
\bottomrule
\end{tabular}
\end{table}

A broader historical comparison, evaluating eleven representative
checkpoints spanning the entire project history on \emph{both} the
grasp0 and grasp10 CMC eval sets (whereas nearly all earlier reported
numbers, including Table~\ref{tab:top10}, are grasp0-only), confirms
v102 and v152 remain jointly the two strongest configurations discovered
so far, essentially tied (282.85 vs.\ 282.89 combined sum), with the
IRLS-robustified variant (v151) close behind but not ahead -- consistent
with Section~\ref{sec:bg_removal}'s finding that IRLS does not fully
resolve the underlying contamination issue. Full per-checkpoint,
per-dataset results are provided in the project repository.

\section{Discussion: Diagnosing the Current Ceiling}
\label{sec:discussion}

Two structural factors, not any single implementation detail, explain why
this project's $\sim$150 experiments cluster so tightly around the same
performance band. First, the only external supervisory signal is optical
flow from a flow estimator trained on synthetic data and applied,
zero-shot, to a visually very different domain (laparoscopic video with
specular highlights, blood, smoke, and fast, thin-object motion); a prior
diagnosis in this project directly measured that the network's own
mask-prediction error correlates strongly with RAFT's own forward-backward
cycle-consistency confidence, i.e.\ a substantial share of the error
budget is attributable to the supervisory signal itself, not to anything
downstream of it. An attempt to remove RAFT from the loop entirely
(replacing it with a self-taught, in-network correlation-based flow)
collapsed via self-referential feedback and was abandoned -- to date,
every subsequent experiment, including this paper's own, has continued to
regress against the same fixed-quality external signal rather than
improving it at the source.

Second, the instrument occupies roughly 1\% of frame area, and RAFT's
own per-pixel error does not shrink with object size -- so the
effective signal-to-noise ratio available to the instrument channel is
structurally worse than for tissue or background, which can average
flow-estimation noise over a much larger area. This offers a unifying
explanation for a pattern that recurs across many, otherwise unrelated
experiments in this project's history: tissue mIoU improves steadily over
training while instrument mIoU oscillates or degrades. We verified this
is not unique to the background-removal mechanism diagnosed in
Section~\ref{sec:bg_removal} -- a milder version of the same
instrument-declines/tissue-improves tension is visible even in the
project's own historical champion configuration.

Neither factor is something further architecture search, on this
project's current data scale, is likely to resolve on its own -- both
Section~\ref{sec:archsearch} (architecture) and
Section~\ref{sec:dino_graph} (adding a new, independent information
source) were tested directly and, at best, matched rather than clearly
exceeded the two signal-quality corrections of
Section~\ref{sec:signal_quality}.

\section{Future Work}
\label{sec:future}

We highlight one design principle, surfaced repeatedly and independently
across this project's history, as the guiding constraint for future
work: any attempt to denoise this architecture's already-weak supervisory
signal must be \emph{filtering}, not \emph{smoothing} -- it must
selectively identify and remove genuine noise while leaving the
underlying (already scarce) signal untouched, rather than uniformly
blurring/averaging in a way that degrades signal and noise together.
Every one of this project's clean wins (Section~\ref{sec:signal_quality})
fits this description; every negative or inconclusive result we can point
to a clear cause for (Section~\ref{sec:bg_removal}'s IRLS,
Section~\ref{sec:sliding_window}'s taper/nearest-center alternatives) is,
in one way or another, a smoothing operation applied where a filtering
operation was needed. We note, separately, that this architecture's own
residual-correction and affine-relaxation machinery (Section
\ref{sec:base_arch}) already gives it real tolerance for imperfect rigid
motion assumptions by design -- the failure mode we diagnose is
specifically an unfiltered contamination of the \emph{background}
estimate, not a lack of robustness in the model itself.

Concretely, we propose two directions motivated by this constraint and by
\cite{batrack2025}'s ablation (Section~\ref{sec:related_work}):

\begin{itemize}
\item \textbf{A dedicated, separately-trained motion-decoupling
sub-network.} \cite{batrack2025} finds that a single network cannot
reliably learn both total motion and the dynamic-only component
simultaneously, and that a small, dedicated network for the dynamic
component -- rather than reusing the main segmentation branch -- is
required. This project's own attempted fix in Section~\ref{sec:bg_removal}
reused the \emph{same} segmentation branch (an EMA copy of it) for both
jobs; a genuinely independent, purpose-built sub-network whose only task
is instrument/background discrimination for the background-motion fit,
supervised indirectly through the effect it has on the fit's own
downstream quality, is a natural next step and has not yet been tried.
\item \textbf{Cross-frame evidence aggregation.} All current inference is
per-frame independent; RAFT's own errors are therefore never
cross-checked against any other estimate. This project has already
computed, but never used for training, dense multi-gap flow across up to
28 frame-pair combinations per continuous-video case -- a natural,
already-available resource for aggregating multiple independent flow
estimates of the same underlying motion into a single, lower-noise
target, rather than trusting any single-frame-pair RAFT estimate
outright.
\end{itemize}

Lower-priority directions surfaced during this project but not yet
pursued include self-supervised fine-tuning of the flow estimator itself
on the target domain (photometric-consistency-based, requiring no labels)
to address the domain-gap half of Section~\ref{sec:discussion}'s
diagnosis directly, and multi-scale test-time ensembling as a cheaper,
no-retraining alternative to Section~\ref{sec:this_session}'s multi-scale
training fix.

\section*{Acknowledgment}
% fill in supervisor / lab / compute acknowledgments here

\bibliographystyle{IEEEtran}
\bibliography{refs}

\end{document}
